\subsection{Conceptual Position within Modified-Gravity and Inhomogeneous Frameworks}
  \label{subsec:conceptual_position}

  The phenomenology discussed in this work exhibits partial overlap
  with acceleration-based approaches such as Modified Newtonian Dynamics (MOND),
  and the cosmological discussion involves large-scale underdensities.
  It is therefore important to clarify the conceptual position of the
  present framework within this broader landscape.

  \paragraph{Relation to MOND-like phenomenology.}

    At galactic accelerations below $a_\star$,
    Eq.~\eqref{eq:modified_poisson} becomes formally equivalent
    to quasi-linear MOND in the weak-field regime.
    Both approaches identify a characteristic acceleration scale
    below which Newtonian expectations break down.

    The distinction is structural rather than mathematical.
    In MOND, the acceleration scale $a_0$ is introduced as a
    fundamental modification of the gravitational law.
    In the present framework, the same scaling behavior arises
    as the low-density limit of a bounded-response principle.
    The scale $a_\star$ characterizes the onset of saturation in
    the effective description, not a new force law.

    Crucially, the extension to cosmological inference does not
    follow from standard MOND formulations.
    Here, both galactic dynamics and the environment-dependent
    Hubble offset derive from the same empirically calibrated scale
    $a_\star$, linking two regimes that are usually treated independently.

  \paragraph{Distinction from inhomogeneous cosmologies.}

    Models based on large-scale metric inhomogeneities,
    such as Lemaître--Tolman--Bondi (LTB) constructions,
    attempt to account for late-time discrepancies by modifying
    the background geometry itself.
    In such scenarios, the Hubble parameter becomes position dependent
    because the spacetime metric deviates from the homogeneous
    Friedmann--Lemaître--Robertson--Walker solution.

    The present framework does not modify the background metric
    and does not introduce a position-dependent Friedmann solution.
    The global expansion history remains that of standard cosmology.
    The predicted offset arises instead from an environment-dependent
    modification of the effective kinematic inference once local
    accelerations fall below the saturation scale $a_\star$.

    The mechanism therefore does not require a special observer
    location or a finely tuned void profile.
    The shift in the locally inferred expansion rate scales with
    the surrounding density contrast according to
    Eq.~\eqref{eq:hloc}, while preserving global homogeneity
    at the level of the background solution.

  \paragraph{Structural distinction.}

    The unifying feature of the framework is the existence of a
    universal bound on effective response.
    This bound is not environment dependent.
    What varies across environments is the density of relations
    contributing to the effective projection.

    Dense systems remain in the unsaturated regime,
    while extended low-density regions approach saturation.
    The resulting correlation between galactic kinematics and
    late-time Hubble inference is therefore not a reparameterization
    of dark sector physics nor a modification of the cosmological metric,
    but a manifestation of a single bounded-response structure
    operating across scales.
